\ifx\allfiles\undefined
\documentclass[12pt, a4paper, oneside, UTF8]{ctexbook}
\def\path{../config}
\input{\path/package.tex}

% 定理环境设置
\input{\path/theorem.tex}

\input{\path/custom.tex}

\input{\path/cover_package} 

% 封面页设置
\def\myTitle{量子力学：理论与应用
}
\def\mySubtitle{Quantum Mechanics: Theory and Application}
\def\myAuthor{Max}
\def\myCreateDateCover{2026年10月11日}
\def\myUpdateDateCover{ \today}
\def\myDateForeword{}
% 前言页设置
\def\myForeword{前言}
\def\myForewordText{

}
\def\mySubheading{}
\begin{document}
% \input{\path/cover_text.tex}

%%%%%%%% 前言页 %%%%%%%%

\newpage
\thispagestyle{empty}
\begin{center}
    \Huge\textbf{\myForeword}
\end{center}
\myForewordText
\begin{flushright}
    \begin{tabular}{c}
        \myDateForeword
    \end{tabular}
\end{flushright}

% 目录页
\newpage
\pagestyle{plain}
\setcounter{page}{1}
\pagenumbering{Roman}
\tableofcontents

% 正文页
\newpage
\pagenumbering{arabic}
\setcounter{chapter}{-1}
\setcounter{page}{1}

\pagestyle{fancy}
% \fancyhead[L]{量子力学笔记}
% \fancyhead[C]{中间页眉}
% \fancyhead[R]{右侧页眉}
% \fancyfoot[C]{\thepage}
\fancyhf{}
% 设置页眉
\fancyhead[L]{\leftmark}    % 左侧：文章标题（\leftmark 通常是章标题）
\fancyhead[C]{ }                       % 中间：留空（或自定义内容）
\fancyhead[R]{\rightmark}   % 右侧：节标题（\rightmark 通常是节标题）


% 设置页脚（示例：中间显示页码）
\fancyfoot[C]{\thepage}

\renewcommand{\headrulewidth}{0.4pt}
\renewcommand{\footrulewidth}{0pt}








\else
\fi

\chapter{中心力场}
\begin{examct}{}{}
    中心力场中粒子运动的性质、轨道角动量、球方势阱、氢原子
\end{examct}
\begin{examrqm}{}{}
    \begin{enumerate}
        \item 掌握中心力场中粒子运动的性质，会求解轨道角动量的本征值和本征函数，掌握球方势阱的基本性质，会求解氢原子的本征值和本征函数
        \item 能综合运用所学知识解决实际问题
    \end{enumerate}
\end{examrqm}

\section{中心力场的一般理论}
\subsection{中心力场的定态薛定谔方程}
球坐标下体系哈密顿量为:
\begin{equation}
    \hat{H} = -\frac{\hbar^2}{2m}\frac{1}{r}\frac{\partial}{\partial r} (r^2\frac{\partial}{\partial r})+ \frac{\hat{L}^2}{2mr^2} + V(r) 
\end{equation}

定态薛定谔方程：
\begin{equation}
\left[ -\frac{\hbar^2}{2m}\frac{1}{r}\frac{\partial}{\partial r} (r^2\frac{\partial}{\partial r})+ \frac{\hat{L}^2}{2mr^2} + V(r) \right] \psi(r,\theta,\phi) = E\psi(r,\theta,\phi)
\end{equation}
\subsection{中心力场的定态波函数}
波函数分为径向部分和角向部分:
\begin{equation}
    \psi(r,\theta,\phi) = Y_{l,m}(\theta,\phi)R(r)
\end{equation}

径向方程：
\begin{equation}
    \left[ -\frac{\hbar^2}{2m}\frac{1}{r}\frac{\partial}{\partial r} (r^2\frac{\partial}{\partial r}) + \frac{l(l+1)\hbar^2}{2mr^2} + V(r) \right] R(r) = ER(r)
\end{equation}

约化径向方程：
令$u(r) = R(r)$
\begin{equation}
    \left[ -\frac{\hbar^2}{2m}\frac{\partial^2}{\partial r^2} + \frac{l(l+1)\hbar^2}{2mr^2} + V(r) \right] u(r) = Eu(r)
\end{equation}



\section{轨道角动量的本征值和本征函数}


\section{球方势阱的基本性质}
\section{氢原子的本征值和本征函数}



\ifx\allfiles\undefined
\end{document}
\fi